\Page{Chapitre 01\quad /\quad Ondes et signaux}{Formation des images}{chapitre-01}
\vspace{2mm}
{\otLight\fontsize{18}{23}\selectfont\color{otNavy}Optique géométrique, instruments et guidage}\par
\vspace{3mm}
Comment une lentille peut-elle former une image sur un écran, donner une image virtuelle ou guider la lumière vers un capteur ? Le même modèle de rayons permet d'étudier ces situations, à condition de connaître ses hypothèses et ses limites.
\Sub{Les objectifs du chapitre}
\begin{itemize}
\item Caractériser une source lumineuse et relier fréquence, longueur d'onde, vitesse et indice.
\item Appliquer les lois de Snell-Descartes, établir un seuil de réflexion totale et tracer le trajet d'un rayon.
\item Construire et calculer des images par un miroir plan ou une lentille mince, avec des distances algébriques cohérentes.
\item Relier la netteté à un détecteur et comprendre les réglages d'un œil ou d'un appareil photographique.
\item Établir les conditions de guidage et l'étalement temporel d'une fibre à saut d'indice.
\item Former une image au laboratoire, exploiter des mesures et discuter une incertitude.
\end{itemize}
\Sub{Les acquis à réactiver}
Conversions d'unités, trigonométrie du triangle rectangle, angles en radians, proportionnalité et résolution d'équations simples. Le diagnostic p.\,\pageref{diagnostic} permet de faire le point avant de commencer.
\begin{loi}{Le parcours du chapitre}
\begin{tabularx}{\linewidth}{@{}Y r@{}}
\tocrow{Notions et capacités du programme}{concordance}
\tocrow{Cours et schémas}{c-sources}
\tocrow{TP et activité numérique}{c-tp}
\tocrow{24 exercices progressifs}{td1}
\tocrow{3 défis, DM et DS}{defis1}
\tocrow{Indications puis corrigés}{indices}
\tocrow{Parcours sur des sujets authentiques}{authentiques}
\end{tabularx}
\end{loi}
